%----------------------------------------------------------
% 제6장. AI 거버넌스 도구 및 기술
%----------------------------------------------------------

\chapter{AI 거버넌스 도구 및 기술}
\label{ch:governance_tools}

AI 거버넌스의 효과적인 실행을 위해서는 적절한 Tool와 기술 인프라가 필수적이다. 
수작업으로 Fairness을 평가하고, 수동으로 모델 카드\cite{mitchell2019model}를 작성하며, 
스프레드시트로 AI 시스템을 관리하는 것은 소수의 모델에서는 가능하지만, 
조직의 AI 활용이 확대됨에 따라 지속 가능하지 않다.

본 장에서는 AI 거버넌스를 기술적으로 Support하는 도구 생태계를 소개하고, 
이를 조직의 MLOps\cite{kreuzberger2023machine} Pipeline에 통합하는 방법을 다룬다. 
오픈소스 Tool와 상용 Platform을 균형 있게 다루며, 
각 Tool의 장단점과 적용 시 고려사항을 분석한다.


%----------------------------------------------------------
\section{AI 거버넌스 기술 스택 개요}
\label{sec:6.1}

\subsection{거버넌스 기술 스택 아키텍처}
\label{sec:6.1.1}

AI 거버넌스 기술 스택은 데이터 수집부터 모델 Retire까지 
전체 AI 생명주기를 Support해야 한다.

\begin{figure}[htbp]
\centering
\resizebox{\textwidth}{!}{%
\begin{tikzpicture}[
  scale=0.85, transform shape,
  layer/.style={rectangle, draw, rounded corners, minimum width=14cm, minimum height=1.2cm, align=center, font=\small},
  sublayer/.style={rectangle, draw, minimum width=4cm, minimum height=0.8cm, align=center, font=\scriptsize},
]

% Layer들
\node[layer, fill=blue!10] (l1) at (0, 6) {\textbf{거버넌스 대시보드 및 보고}\\경영진 대시보드, 규제 Report, 감사 Tracking};

\node[layer, fill=green!10] (l2) at (0, 4.5) {\textbf{정책 및 Workflow 관리}\\승인 Gate, RACI 관리, 문서 템플릿, 체크리스트};

\node[layer, fill=yellow!10] (l3) at (0, 3) {\textbf{평가 및 테스트 도구}\\공정성 평가\cite{bellamy2019ai}, 설명가능한 AI\cite{ribeiro2016why}(설명가능한 AI(XAI)), 견고성 테스트, Privacy 검증};

\node[layer, fill=orange!10] (l4) at (0, 1.5) {\textbf{모니터링\cite{klaise2021monitoring} 및 관측성\cite{shankar2022operationalizing}}\\성능 모니터링, 드리프트\cite{paleyes2022challenges} 감지, 알림, Logging};

\node[layer, fill=red!10] (l5) at (0, 0) {\textbf{MLOps based 인프라}\\Model 레지스트리, 피처 스토어, 실험 추적, CI/CD\cite{sato2019continuous}};

% 수직 통합 요소
\node[rectangle, draw, fill=purple!10, minimum width=1.5cm, minimum height=6cm, align=center, font=\scriptsize, rotate=90]
 at (8, 3) {AI 시스템 레지스트리};

\node[rectangle, draw, fill=gray!10, minimum width=1.5cm, minimum height=6cm, align=center, font=\scriptsize, rotate=90]
 at (-8, 3) {감사 로그 및 버전 관리};

\end{tikzpicture}
}
\caption{AI 거버넌스 기술 스택 아키텍처}
\label{fig:tech_stack_architecture}
\end{figure}


\subsection{도구 카테고리별 개요}
\label{sec:6.1.2}

\begin{table}[htbp]
\centering
\caption{AI 거버넌스 도구 카테고리 및 대표 도구}
\label{tab:tool_categories}
\begin{tabularx}{\textwidth}{p{2.5cm}p{4cm}X}
\toprule
\textbf{카테고리} & \textbf{오픈소스 도구} & \textbf{상용 플랫폼} \\
\midrule
공정성 평가 & 
Fairlearn\cite{mehrabi2021survey}, AIF360\cite{saleiro2018aequitas}, Aequitas, Themis-ML & 
Google What-If 도구, AWS SageMaker Clarify, Azure Fairlearn \\
\midrule
설명가능성\cite{arrieta2020explainable} (설명가능한 AI(설명가능한 AI(XAI))) & 
SHAP\cite{lundberg2017unified}, LIME, Captum\cite{sokol2020explainability}, InterpretML, Alibi & 
Google Explainable AI, AWS SageMaker Debugger \\
\midrule
모델 레지스트리 & 
MLflow\cite{zaharia2018accelerating}, DVC, ModelDB & 
AWS SageMaker Model 레지스트리, Azure ML, Weights \& Biases \\
\midrule
모니터링 & 
Evidently\cite{sculley2015hidden} AI, NannyML, Alibi Detect, WhyLabs & 
Arize AI, Fiddler AI, Arthur AI, Aporia \\
\midrule
Privacy & 
Opacus, diffprivlib, PySyft, TensorFlow Privacy & 
Google Privacy Sandbox, AWS Clean Rooms \\
\midrule
문서화\cite{arnold2019factsheets} & 
Model Card\cite{gebru2021datasheets} Toolkit, FactSheets360 & 
Credo AI, Holistic AI \\
\midrule
통합된 플랫폼 & 
- & 
IBM AI 거버넌스, Credo AI, Holistic AI, Monitaur \\
\bottomrule
\end{tabularx}
\end{table}


\subsection{도구 선정 기준}
\label{sec:6.1.3}

조직에 적합한 AI 거버넌스 Tool를 선정할 때 고려해야 할 기준을 제시한다.

\begin{table}[htbp]
\centering
\caption{AI 거버넌스 도구 선정 기준}
\label{tab:tool_selection_criteria}
\begin{tabularx}{\textwidth}{p{3cm}X}
\toprule
\textbf{기준} & \textbf{평가 항목} \\
\midrule
기능 완전성 & 
필요한 거버넌스 기능(공정성, 설명가능한 AI(설명가능한 AI(XAI)), 모니터링 등) Support 여부; 
확장 가능성 \\
\midrule
기존 시스템 통합 & 
현재 사용 중인 ML 프레임워크(sklearn, PyTorch, TensorFlow), 
MLOps 도구(MLflow, Kubeflow), 
클라우드 플랫폼(AWS, GCP, Azure)과의 호환성 \\
\midrule
확장성 & 
대규모 데이터셋, 다수의 모델, 높은 트래픽 처리 능력 \\
\midrule
사용 용이성 & 
학습 곡선, 문서화 품질, 커뮤니티 Support, UI/UX \\
\midrule
규제 준수 Support & 
시행령, EU AI Act 등 규제 요구사항 템플릿 제공; 
감사 추적 기능 \\
\midrule
비용 & 
라이선스 비용, 인프라 비용, 운영 비용, TCO \\
\midrule
보안 및 Privacy & 
데이터 보안, 접근 제어, 온프레미스 배포 옵션 \\
\midrule
벤더 안정성 & 
오픈소스 커뮤니티 활성도; 상용 벤더의 재무 안정성, 로드맵 \\
\bottomrule
\end{tabularx}
\end{table}

%----------------------------------------------------------
\section{공정성 평가 및 편향 완화 도구}
\label{sec:6.2}

\subsection{Fairlearn 심화 활용}
\label{sec:6.2.1}

Fairlearn은 Microsoft가 개발한 오픈소스 공정성 평가 및 완화 라이브러리로, 
본 책에서 주요 Tool로 활용하고 있다. 
2025년 업데이트된 0.14.0 버전의 주요 기능을 상세히 다룬다.

\begin{lstlisting}[language=Python, caption={Fairlearn 0.14.0 Advanced Usage}, label={lst:fairlearn_advanced}]
"""
Fairlearn 0.14.0 Advanced Usage Guide

- Multi-dimensional fairness analysis using MetricFrame
- Intersectionality analysis for multiple sensitive attributes
- Defining custom metrics
- Comparing bias mitigation algorithms
"""

import numpy as np
import pandas as pd
from sklearn.model_selection import train_test_split
from sklearn.ensemble import GradientBoostingClassifier
from sklearn.metrics import (
 accuracy_score, precision_score, recall_score, 
 f1_score, roc_auc_score
)

from fairlearn.metrics import (
 MetricFrame,
 demographic_parity_difference,
 demographic_parity_ratio,
 equalized_odds_difference,
 selection_rate,
 true_positive_rate,
 false_positive_rate,
 count,
)
from fairlearn.reductions import (
 ExponentiatedGradient,
 GridSearch,
 DemographicParity,
 EqualizedOdds,
 TruePositiveRateParity,
 ErrorRateParity,
)
from fairlearn.postprocessing import ThresholdOptimizer


class AdvancedFairnessAnalyzer:
 """
 Advanced fairness analyzer based on Fairlearn 0.14.0
 
 Supports multiple sensitive attributes, intersectional analysis, and custom metrics
 """
 
 def __init__(self):
 # Base metric set
 self.base_metrics = {
 'accuracy': accuracy_score,
 'precision': precision_score,
 'recall': recall_score,
 'f1': f1_score,
 'selection_rate': selection_rate,
 'tpr': true_positive_rate,
 'fpr': false_positive_rate,
 'count': count,
 }
 
 def create_intersectional_groups(
 self,
 sensitive_features: pd.DataFrame,
 columns: list,
 ) -> pd.Series:
 """
 Create intersectional groups
 
 Example: gender x age_group -> 'female_young', 'male_senior', etc.
 """
 return sensitive_features[columns].apply(
 lambda row: '_'.join(str(v) for v in row), axis=1
 )
 
 def comprehensive_fairness_analysis(
 self,
 y_true: np.ndarray,
 y_pred: np.ndarray,
 sensitive_features: pd.DataFrame,
 intersectional_columns: list = None,
 ) -> dict:
 """
 Comprehensive fairness analysis
 
 Performs analysis for individual sensitive attributes + intersectional groups
 """
 
 results = {
 'individual_attributes': {},
 'intersectional': {},
 'summary': {},
 }
 
 # 1. Individual sensitive attribute analysis
 for col in sensitive_features.columns:
 mf = MetricFrame(
 metrics=self.base_metrics,
 y_true=y_true,
 y_pred=y_pred,
 sensitive_features=sensitive_features[col],
 )
 
 results['individual_attributes'][col] = {
 'by_group': mf.by_group.to_dict(),
 'overall': mf.overall.to_dict(),
 'difference': mf.difference().to_dict(),
 'ratio': mf.ratio().to_dict(),
 'group_min': mf.group_min().to_dict(),
 'group_max': mf.group_max().to_dict(),
 }
 
 # 2. Intersectional analysis (if specified)
 if intersectional_columns and len(intersectional_columns) >= 2:
 intersectional_groups = self.create_intersectional_groups(
 sensitive_features, intersectional_columns
 )
 
 mf_inter = MetricFrame(
 metrics=self.base_metrics,
 y_true=y_true,
 y_pred=y_pred,
 sensitive_features=intersectional_groups,
 )
 
 results['intersectional'] = {
 'columns': intersectional_columns,
 'by_group': mf_inter.by_group.to_dict(),
 'difference': mf_inter.difference().to_dict(),
 'worst_group': mf_inter.group_min().to_dict(),
 }
 
 # 3. Summary metrics
 results['summary'] = {
 'dp_difference': demographic_parity_difference(
 y_true, y_pred,
 sensitive_features=sensitive_features[sensitive_features.columns[0]]
 ),
 'dp_ratio': demographic_parity_ratio(
 y_true, y_pred,
 sensitive_features=sensitive_features[sensitive_features.columns[0]]
 ),
 'eo_difference': equalized_odds_difference(
 y_true, y_pred,
 sensitive_features=sensitive_features[sensitive_features.columns[0]]
 ),
 }
 
 return results
 
 def compare_mitigation_algorithms(
 self,
 X_train: np.ndarray,
 y_train: np.ndarray,
 sensitive_train: np.ndarray,
 X_test: np.ndarray,
 y_test: np.ndarray,
 sensitive_test: np.ndarray,
 base_estimator,
 constraints: list = None,
 ) -> pd.DataFrame:
 """
 Bias mitigation algorithm comparison
 
 Compare multiple constraints and algorithms to find 
 the optimal fairness-accuracy trade-off
 """
 
 if constraints is None:
 constraints = [
 ('Demographic Parity', DemographicParity()),
 ('Equalized Odds', EqualizedOdds()),
 ('TPR Parity', TruePositiveRateParity()),
 ('Error Rate Parity', ErrorRateParity()),
 ]
 
 results = []
 
 # 1. Baseline (No Mitigation)
 base_model = base_estimator.__class__(base_estimator.get_params())
 base_model.fit(X_train, y_train)
 y_pred_base = base_model.predict(X_test)
 
 results.append({
 'method': 'Baseline (No Mitigation)',
 'accuracy': accuracy_score(y_test, y_pred_base),
 'dp_ratio': demographic_parity_ratio(
 y_test, y_pred_base, sensitive_features=sensitive_test
 ),
 'dp_diff': demographic_parity_difference(
 y_test, y_pred_base, sensitive_features=sensitive_test
 ),
 'eo_diff': equalized_odds_difference(
 y_test, y_pred_base, sensitive_features=sensitive_test
 ),
 })
 
 # 2. Apply ExponentiatedGradient for each constraint
 for name, constraint in constraints:
 try:
 mitigator = ExponentiatedGradient(
 estimator=base_estimator.__class__(base_estimator.get_params()),
 constraints=constraint,
 eps=0.01,
 )
 mitigator.fit(X_train, y_train, sensitive_features=sensitive_train)
 y_pred = mitigator.predict(X_test)
 
 results.append({
 'method': f'ExpGrad + {name}',
 'accuracy': accuracy_score(y_test, y_pred),
 'dp_ratio': demographic_parity_ratio(
 y_test, y_pred, sensitive_features=sensitive_test
 ),
 'dp_diff': demographic_parity_difference(
 y_test, y_pred, sensitive_features=sensitive_test
 ),
 'eo_diff': equalized_odds_difference(
 y_test, y_pred, sensitive_features=sensitive_test
 ),
 })
 except Exception as e:
 results.append({
 'method': f'ExpGrad + {name}',
 'accuracy': None,
 'dp_ratio': None,
 'dp_diff': None,
 'eo_diff': None,
 'error': str(e),
 })
 
 # 3. ThresholdOptimizer (Post-processing)
 try:
 threshold_opt = ThresholdOptimizer(
 estimator=base_model,
 constraints="demographic_parity",
 objective="accuracy_score",
 prefit=True,
 )
 threshold_opt.fit(X_train, y_train, sensitive_features=sensitive_train)
 y_pred_thresh = threshold_opt.predict(X_test, sensitive_features=sensitive_test)
 
 results.append({
 'method': 'ThresholdOptimizer (DP)',
 'accuracy': accuracy_score(y_test, y_pred_thresh),
 'dp_ratio': demographic_parity_ratio(
 y_test, y_pred_thresh, sensitive_features=sensitive_test
 ),
 'dp_diff': demographic_parity_difference(
 y_test, y_pred_thresh, sensitive_features=sensitive_test
 ),
 'eo_diff': equalized_odds_difference(
 y_test, y_pred_thresh, sensitive_features=sensitive_test
 ),
 })
 except Exception as e:
 results.append({
 'method': 'ThresholdOptimizer (DP)',
 'error': str(e),
 })
 
 return pd.DataFrame(results)
 
 def generate_fairness_report_latex(
 self,
 analysis_results: dict,
 model_name: str,
 ) -> str:
 """Fairness Analysis LaTeX Report Generation"""
 
 latex = r"""
\begin{table}[htbp]
\centering
\caption{""" + model_name + r""" Fairness Analysis Results}
\begin{tabularx}{\textwidth}{p{3cm}XXXXX}
\toprule
\textbf{Sensitive Attribute} & \textbf{DP Ratio} & \textbf{DP Diff} & \textbf{EO Diff} & \textbf{4/5 Rule} \\
\midrule
"""
 
 for attr, result in analysis_results.get('individual_attributes', {}).items():
 dp_ratio = result.get('ratio', {}).get('selection_rate', 0)
 dp_diff = result.get('difference', {}).get('selection_rate', 0)
 
 # EO diff is the maximum of TPR and FPR differences
 tpr_diff = result.get('difference', {}).get('tpr', 0)
 fpr_diff = result.get('difference', {}).get('fpr', 0)
 eo_diff = max(abs(tpr_diff) if tpr_diff else 0, 
 abs(fpr_diff) if fpr_diff else 0)
 
 passes = '[PASS]' if dp_ratio and dp_ratio >= 0.8 else '[FAIL]'
 
 latex += f"{attr} & {dp_ratio:.4f} & {dp_diff:.4f} & {eo_diff:.4f} & {passes} \\\\\n"
 
 latex += r"""
\bottomrule
\end{tabularx}
\end{table}
"""
 
 return latex
\end{lstlisting}

\subsection{AIF360 (AI 공정성 360)}
\label{sec:6.2.2}

IBM이 개발한 AIF360은 Fairlearn과 함께 가장 널리 사용되는 공정성 Tool이다. 
Fairlearn이 sklearn 친화적이라면, AIF360은 자체 데이터셋 클래스와 
더 다양한 메트릭을 제공한다.

\begin{table}[htbp]
\centering
\caption{Fairlearn vs AIF360 Comparison}
\label{tab:fairlearn_vs_aif360}
\begin{tabularx}{\textwidth}{p{3cm}XX}
\toprule
\textbf{항목} & \textbf{Fairlearn 0.14.0} & \textbf{AIF360 0.6.1} \\
\midrule
개발사 & Microsoft & IBM \\
\midrule
인터페이스 & sklearn 호환, pandas 친화적 & 자체 Dataset 클래스 \\
\midrule
메트릭 수 & 약 15개 & 약 70개 \\
\midrule
완화 알고리즘 & ExponentiatedGradient, GridSearch, ThresholdOptimizer & Reweighing, Prejudice Remover, Calibrated EqOdds 등 20개+ \\
\midrule
시각화 & 기본 제공 & FairnessDashboard \\
\midrule
장점 & sklearn 통합 용이, 직관적 API & 풍부한 메트릭, 다양한 알고리즘 \\
\midrule
단점 & 상대적으로 적은 알고리즘 & 학습 곡선 높음, 무거움 \\
\midrule
권장 사용 & sklearn based 파이프라인, 빠른 적용 & 연구, 심층 분석, 다양한 완화 기법 비교 \\
\bottomrule
\end{tabularx}
\end{table}


%----------------------------------------------------------
\section{설명가능성(설명가능한 AI(설명가능한 AI(XAI))) 도구}
\label{sec:6.3}

\subsection{SHAP 심화 활용}
\label{sec:6.3.1}

SHAP(SHapley Additive exPlanations)은 가장 널리 사용되는 설명가능한 AI(설명가능한 AI(XAI)) 라이브러리이다. 
0.46.0 버전의 주요 기능과 실무 활용법을 다룬다.

\begin{lstlisting}[language=Python, caption={SHAP 0.46.0 Advanced Usage}, label={lst:shap_advanced}]
"""
SHAP 0.46.0 Advanced Usage Guide

- Selecting optimal Explainer by model type
- Large-scale data processing
- Customizing visualizations
- Governance report integration
"""

import numpy as np
import pandas as pd
import shap
from typing import Optional, Dict, List, Union, Any
import warnings

class SHAPGovernanceExplainer:
 """
 SHAP-based governance explanation generator
 
 Optimized for model types, handles large-scale data,
 and supports governance report formats
 """
 
 # Model-Explainer 
 EXPLAINER_MAP = {
 'tree': shap.TreeExplainer,
 'linear': shap.LinearExplainer,
 'deep': shap.DeepExplainer,
 'gradient': shap.GradientExplainer,
 'kernel': shap.KernelExplainer,
 }
 
 # sklearn based Model
 TREE_MODELS = [
 'RandomForestClassifier', 'RandomForestRegressor',
 'GradientBoostingClassifier', 'GradientBoostingRegressor',
 'XGBClassifier', 'XGBRegressor',
 'LGBMClassifier', 'LGBMRegressor',
 'CatBoostClassifier', 'CatBoostRegressor',
 'DecisionTreeClassifier', 'DecisionTreeRegressor',
 'ExtraTreesClassifier', 'ExtraTreesRegressor',
 ]
 
 # Model
 LINEAR_MODELS = [
 'LogisticRegression', 'LinearRegression',
 'Ridge', 'Lasso', 'ElasticNet',
 'SGDClassifier', 'SGDRegressor',
 ]
 
 def __init__(
 self,
 model: Any,
 feature_names: Optional[List[str]] = None,
 max_samples_for_kernel: int = 100,
 ):
 """
 Parameters
 ----------
 model : Any
 Model to explain
 feature_names : list, optional
 List of feature names
 max_samples_for_kernel : int
 Max background samples for KernelExplainer
 """
 self.model = model
 self.feature_names = feature_names
 self.max_samples_for_kernel = max_samples_for_kernel
 self.explainer = None
 self.explainer_type = None
 self.background_data = None
 
 def _detect_model_type(self) -> str:
 """Auto-detect model type"""
 
 model_class = type(self.model).__name__
 
 if model_class in self.TREE_MODELS:
 return 'tree'
 elif model_class in self.LINEAR_MODELS:
 return 'linear'
 elif 'keras' in str(type(self.model)).lower():
 return 'deep'
 elif 'torch' in str(type(self.model)).lower():
 return 'gradient'
 else:
 return 'kernel'
 
 def fit(
 self,
 X_background: Union[np.ndarray, pd.DataFrame],
 model_type: Optional[str] = None,
 ) -> 'SHAPGovernanceExplainer':
 """
 Initialize Explainer
 
 Parameters
 ----------
 X_background : array-like
 Background data (used in KernelExplainer, DeepExplainer)
 model_type : str, optional
 One of 'tree', 'linear', 'deep', 'gradient', 'kernel'.
 Auto-detected if None.
 """
 
 # Model self.explainer_type = model_type or self._detect_model_type()
 
 # Data 
 if isinstance(X_background, pd.DataFrame):
 self.background_data = X_background.values
 if self.feature_names is None:
 self.feature_names = X_background.columns.tolist()
 else:
 self.background_data = X_background
 
 # Explainer 
 if self.explainer_type == 'tree':
 # 0.46.0: check_additivity self.explainer = shap.TreeExplainer(
 self.model,
 feature_perturbation='interventional',
 )
 
 elif self.explainer_type == 'linear':
 self.explainer = shap.LinearExplainer(
 self.model,
 self.background_data,
 )
 
 elif self.explainer_type in ['deep', 'gradient']:
 # Model: Number of samples 
 bg_sample = self.background_data
 if len(bg_sample) > self.max_samples_for_kernel:
 indices = np.random.choice(
 len(bg_sample), 
 self.max_samples_for_kernel, 
 replace=False
 )
 bg_sample = bg_sample[indices]
 
 if self.explainer_type == 'deep':
 self.explainer = shap.DeepExplainer(self.model, bg_sample)
 else:
 self.explainer = shap.GradientExplainer(self.model, bg_sample)
 
 else:
 # KernelExplainer ( , )
 bg_sample = self.background_data
 if len(bg_sample) > self.max_samples_for_kernel:
 bg_sample = shap.sample(bg_sample, self.max_samples_for_kernel)
 
 # if hasattr(self.model, 'predict_proba'):
 predict_fn = lambda x: self.model.predict_proba(x)[:, 1]
 else:
 predict_fn = self.model.predict
 
 self.explainer = shap.KernelExplainer(predict_fn, bg_sample)
 
 return self
 
 def explain(
 self,
 X: Union[np.ndarray, pd.DataFrame],
 check_additivity: bool = False,
 ) -> shap.Explanation:
 """
 Calculate SHAP values
 
 Parameters
 ----------
 X : array-like
 Instance to explain
 check_additivity : bool
 Whether to perform additivity check (tree models)
 """
 
 if self.explainer is None:
 raise ValueError("fit() .")
 
 if isinstance(X, pd.DataFrame):
 X = X.values
 
 if self.explainer_type == 'tree':
 shap_values = self.explainer.shap_values(
 X, check_additivity=check_additivity
 )
 else:
 shap_values = self.explainer.shap_values(X)
 
 # Explanation Generation
 if isinstance(shap_values, list):
 # : 1 ( ) 
 shap_values = shap_values[1] if len(shap_values) > 1 else shap_values[0]
 
 return shap.Explanation(
 values=shap_values,
 base_values=self.explainer.expected_value if not isinstance(
 self.explainer.expected_value, list
 ) else self.explainer.expected_value[1],
 data=X,
 feature_names=self.feature_names,
 )
 
 def explain_single(
 self,
 instance: Union[np.ndarray, pd.Series],
 top_k: int = 5,
 ) -> Dict:
 """
 Explain a single instance
 
 Returns
 -------
 dict
 {
 'shap_values': {feature: value},
 'top_positive': [(feature, value), ...],
 'top_negative': [(feature, value), ...],
 'base_value': float,
 'prediction': float,
 }
 """
 
 if isinstance(instance, pd.Series):
 instance = instance.values
 
 instance = instance.reshape(1, -1)
 explanation = self.explain(instance)
 
 shap_dict = {}
 for i, feat in enumerate(self.feature_names or range(len(explanation.values[0]))):
 shap_dict[feat] = explanation.values[0][i]
 
 # Alignment
 sorted_positive = sorted(
 [(k, v) for k, v in shap_dict.items() if v > 0],
 key=lambda x: x[1], reverse=True
 )[:top_k]
 
 sorted_negative = sorted(
 [(k, v) for k, v in shap_dict.items() if v < 0],
 key=lambda x: x[1]
 )[:top_k]
 
 return {
 'shap_values': shap_dict,
 'top_positive': sorted_positive,
 'top_negative': sorted_negative,
 'base_value': float(explanation.base_values),
 'prediction': float(explanation.base_values + explanation.values[0].sum()),
 }
 
 def global_importance(
 self,
 X: Union[np.ndarray, pd.DataFrame],
 method: str = 'mean_abs',
 ) -> pd.DataFrame:
 """
 Calculate global feature importance
 
 Parameters
 ----------
 method : str
 'mean_abs': Mean absolute value
 'mean': Mean (with direction)
 'max_abs': Max absolute value
 """
 
 explanation = self.explain(X)
 
 if method == 'mean_abs':
 importance = np.abs(explanation.values).mean(axis=0)
 elif method == 'mean':
 importance = explanation.values.mean(axis=0)
 elif method == 'max_abs':
 importance = np.abs(explanation.values).max(axis=0)
 else:
 raise ValueError(f"Unknown method: {method}")
 
 df = pd.DataFrame({
 'feature': self.feature_names or [f'feature_{i}' for i in range(len(importance))],
 'importance': importance,
 })
 
 return df.sort_values('importance', ascending=False).reset_index(drop=True)
 
 def generate_governance_explanation(
 self,
 instance: np.ndarray,
 audience: str = 'technical',
 feature_descriptions: Optional[Dict[str, str]] = None,
 ) -> str:
 """
 Generate explanation text for governance
 
 Parameters
 ----------
 audience : str
 'end_user', 'business', 'technical', 'regulator'
 """
 
 result = self.explain_single(instance)
 
 if audience == 'end_user':
 # For end users (Simplified)
 text = "## Key Decision Factors\n\n"
 text += "The following factors influenced the decision:\n\n"
 
 text += "### Impact\n"
 for feat, val in result['top_positive'][:3]:
 desc = feature_descriptions.get(feat, feat) if feature_descriptions else feat
 text += f"- {desc}\n"
 
 if result['top_negative']:
 text += "\n### Impact\n"
 for feat, val in result['top_negative'][:3]:
 desc = feature_descriptions.get(feat, feat) if feature_descriptions else feat
 text += f"- {desc}\n"
 
 elif audience == 'technical':
 # For technical teams (Detailed metrics)
 text = "## SHAP Analysis Results\n\n"
 text += f"- Base Value: {result['base_value']:.4f}\n"
 text += f"- Prediction: {result['prediction']:.4f}\n\n"
 
 text += "### SHAP Values (Top 10)\n\n"
 text += "| Feature | SHAP Value | Direction |\n"
 text += "|---------|------------|-----------|--\n"
 
 sorted_all = sorted(
 result['shap_values'].items(),
 key=lambda x: abs(x[1]), reverse=True
 )[:10]
 
 for feat, val in sorted_all:
 direction = "" if val > 0 else ""
 text += f"| {feat} | {val:.4f} | {direction} |\n"
 
 elif audience == 'regulator':
 # For regulators (Compliance with Enforcement Decree)
 text = "## AI Decision Explanation (Article 23)\n\n"
 text += f"### Prediction Result\n"
 text += f"- Base Value: {result['base_value']:.4f}\n"
 text += f"- Final Prediction: {result['prediction']:.4f}\n\n"
 
 text += "### Major Impact ( 5 )\n\n"
 text += "| | | | Impact |\n"
 text += "|------|------|--------|----------|\n"
 
 sorted_all = sorted(
 result['shap_values'].items(),
 key=lambda x: abs(x[1]), reverse=True
 )[:5]
 
 for i, (feat, val) in enumerate(sorted_all, 1):
 direction = " " if val > 0 else " "
 text += f"| {i} | {feat} | {val:.4f} | {direction} |\n"
 
 text += "\n*This explanation is based on SHAP (SHapley Additive exPlanations) methodology.*\n"
 text += "*Retained for 5 years per Article 15.*\n"
 
 else: # business
 text = "## Business Insights\n\n"
 text += "### Key factors influencing the decision\n\n"
 
 for i, (feat, val) in enumerate(result['top_positive'][:3], 1):
 desc = feature_descriptions.get(feat, feat) if feature_descriptions else feat
 text += f"{i}. {desc} (+{val:.2f})\n"
 
 return text
\end{lstlisting}


\subsection{딥러닝 모델을 위한 설명가능한 AI(설명가능한 AI(XAI)): Captum}
\label{sec:6.3.2}

PyTorch based 딥러닝 모델에는 Captum 라이브러리가 적합하다.

\begin{lstlisting}[language=Python, caption={Captum 0.7.0Advanced Deep Learning with XAI}, label={lst:captum_xai}]
"""
Captum 0.7.0 Usage Guide

- Integrated Gradients
- Layer GradCAM (Vision)
- Attention Visualization
"""

import torch
import torch.nn as nn
import numpy as np
from typing import Optional, Tuple, List
from captum.attr import (
 IntegratedGradients,
 LayerGradCam,
 GuidedGradCam,
 Saliency,
 DeepLift,
 NoiseTunnel,
 visualization as viz,
)

class DeepLearningExplainer:
 """
 Captum-based deep learning model explainer
 """
 
 def __init__(
 self,
 model: nn.Module,
 device: str = 'cpu',
 ):
 self.model = model.to(device)
 self.model.eval()
 self.device = device
 
 def integrated_gradients(
 self,
 inputs: torch.Tensor,
 target: int,
 baseline: Optional[torch.Tensor] = None,
 n_steps: int = 50,
 ) -> torch.Tensor:
 """
 Calculate Integrated Gradients
 
 Reference: Sundararajan et al., "Axiomatic Attribution for Deep Networks" (2017)
 """
 
 ig = IntegratedGradients(self.model)
 
 inputs = inputs.to(self.device)
 if baseline is None:
 baseline = torch.zeros_like(inputs).to(self.device)
 else:
 baseline = baseline.to(self.device)
 
 attributions = ig.attribute(
 inputs,
 baselines=baseline,
 target=target,
 n_steps=n_steps,
 return_convergence_delta=False,
 )
 
 return attributions
 
 def grad_cam(
 self,
 inputs: torch.Tensor,
 target: int,
 layer: nn.Module,
 ) -> torch.Tensor:
 """
 Grad-CAM for CNN
 
 Visualize activation maps for Vision models
 """
 
 grad_cam = LayerGradCam(self.model, layer)
 
 inputs = inputs.to(self.device)
 attributions = grad_cam.attribute(inputs, target=target)
 
 return attributions
 
 def explain_image_classification(
 self,
 image: torch.Tensor,
 target_class: int,
 conv_layer: nn.Module,
 class_names: Optional[List[str]] = None,
 ) -> dict:
 """
 Comprehensive explanation for image classification models
 """
 
 image = image.to(self.device)
 
 # 
 with torch.no_grad():
 output = self.model(image)
 probs = torch.softmax(output, dim=1)
 predicted_class = probs.argmax(dim=1).item()
 confidence = probs[0, predicted_class].item()
 
 # Integrated Gradients
 ig_attr = self.integrated_gradients(image, target_class)
 
 # Grad-CAM
 gradcam_attr = self.grad_cam(image, target_class, conv_layer)
 
 return {
 'predicted_class': predicted_class,
 'predicted_label': class_names[predicted_class] if class_names else str(predicted_class),
 'confidence': confidence,
 'target_class': target_class,
 'integrated_gradients': ig_attr.cpu().numpy(),
 'grad_cam': gradcam_attr.cpu().numpy(),
 }
 
 def explain_tabular(
 self,
 inputs: torch.Tensor,
 target: int,
 feature_names: List[str],
 method: str = 'integrated_gradients',
 ) -> dict:
 """
 Data Model 
 """
 
 inputs = inputs.to(self.device)
 
 if method == 'integrated_gradients':
 ig = IntegratedGradients(self.model)
 attributions = ig.attribute(inputs, target=target)
 elif method == 'deeplift':
 dl = DeepLift(self.model)
 attributions = dl.attribute(inputs, target=target)
 elif method == 'saliency':
 sal = Saliency(self.model)
 attributions = sal.attribute(inputs, target=target)
 else:
 raise ValueError(f"Unknown method: {method}")
 
 attr_values = attributions.cpu().numpy().flatten()
 
 # feature_importance = {
 feat: float(val) 
 for feat, val in zip(feature_names, attr_values)
 }
 
 # Alignment
 sorted_importance = sorted(
 feature_importance.items(),
 key=lambda x: abs(x[1]),
 reverse=True
 )
 
 return {
 'method': method,
 'feature_importance': feature_importance,
 'sorted_importance': sorted_importance,
 'top_positive': [(k, v) for k, v in sorted_importance if v > 0][:5],
 'top_negative': [(k, v) for k, v in sorted_importance if v < 0][:5],
 }
\end{lstlisting}

%----------------------------------------------------------
\section{Model 관리 및 MLOps 통합된}
\label{sec:6.4}

\subsection{MLflow와 거버넌스 통합된}
\label{sec:6.4.1}

MLflow는 가장 널리 사용되는 MLOps Platform으로, 
AI 거버넌스 정보를 통합 관리하는 데 활용할 수 있다.

\begin{lstlisting}[language=Python, caption={MLflow Governance Integration}, label={lst:mlflow_governance}]
"""
MLflow + AI Governance Integration

- Fairness Metric Logging
- Model Card Integration
- Approval Status Management
- Governance Tags
"""

import mlflow
from mlflow.tracking import MlflowClient
from mlflow.models import ModelSignature
from typing import Dict, Optional, List, Any
from datetime import datetime
import json

class MLflowGovernanceIntegration:
 """
 MLflow AI Governance Integrated
 
 Tracking, Model Governance information Integration
 """
 
 # Governance Tags 
 GOVERNANCE_PREFIX = "governance"
 
 # Governance Tags
 GOVERNANCE_TAGS = {
 'risk_level': f'{GOVERNANCE_PREFIX}.risk_level',
 'high_risk_sector': f'{GOVERNANCE_PREFIX}.high_risk_sector',
 'fairness_status': f'{GOVERNANCE_PREFIX}.fairness_status',
 'approval_status': f'{GOVERNANCE_PREFIX}.approval_status',
 'approved_by': f'{GOVERNANCE_PREFIX}.approved_by',
 'approval_date': f'{GOVERNANCE_PREFIX}.approval_date',
 'human_oversight_level': f'{GOVERNANCE_PREFIX}.human_oversight_level',
 'retention_until': f'{GOVERNANCE_PREFIX}.retention_until',
 }
 
 def __init__(
 self,
 tracking_uri: Optional[str] = None,
 experiment_name: Optional[str] = None,
 ):
 if tracking_uri:
 mlflow.set_tracking_uri(tracking_uri)
 
 if experiment_name:
 mlflow.set_experiment(experiment_name)
 
 self.client = MlflowClient()
 
 def log_fairness_metrics(
 self,
 fairness_results: Dict[str, float],
 step: Optional[int] = None,
 ) -> None:
 """Fairness Metric Logging"""
 
 for metric_name, value in fairness_results.items():
 mlflow.log_metric(
 f"fairness.{metric_name}",
 value,
 step=step
 )
 
 # 4/5 dp_ratio = fairness_results.get('dp_ratio', 0)
 passes_4_5 = dp_ratio >= 0.8
 mlflow.log_metric("fairness.passes_4_5_rule", int(passes_4_5))
 
 def log_governance_tags(
 self,
 risk_level: str,
 high_risk_sector: Optional[str] = None,
 fairness_status: str = 'pending',
 human_oversight_level: str = 'HOTL',
 ) -> None:
 """Governance Tags Logging"""
 
 tags = {
 self.GOVERNANCE_TAGS['risk_level']: risk_level,
 self.GOVERNANCE_TAGS['fairness_status']: fairness_status,
 self.GOVERNANCE_TAGS['human_oversight_level']: human_oversight_level,
 }
 
 if high_risk_sector:
 tags[self.GOVERNANCE_TAGS['high_risk_sector']] = high_risk_sector
 
 # 5years Calculation
 retention_until = datetime.now().replace(year=datetime.now().year + 5)
 tags[self.GOVERNANCE_TAGS['retention_until']] = retention_until.isoformat()
 
 mlflow.set_tags(tags)
 
 def log_model_card(
 self,
 model_card: Dict[str, Any],
 artifact_path: str = "governance/model_card.json",
 ) -> None:
 """Model """
 
 # JSON with open("/tmp/model_card.json", "w", encoding="utf-8") as f:
 json.dump(model_card, f, ensure_ascii=False, indent=2, default=str)
 
 mlflow.log_artifact("/tmp/model_card.json", artifact_path="governance")
 
 def log_fairness_report(
 self,
 report_content: str,
 report_format: str = "md",
 ) -> None:
 """Fairness Report """
 
 filename = f"/tmp/fairness_report.{report_format}"
 with open(filename, "w", encoding="utf-8") as f:
 f.write(report_content)
 
 mlflow.log_artifact(filename, artifact_path="governance")
 
 def register_model_with_governance(
 self,
 model_uri: str,
 name: str,
 risk_level: str,
 fairness_metrics: Dict[str, float],
 approval_status: str = 'pending',
 tags: Optional[Dict[str, str]] = None,
 ) -> str:
 """Governance information Model """
 
 # Model 
 result = mlflow.register_model(model_uri, name)
 model_version = result.version
 
 # Governance Tags Setup
 governance_tags = {
 self.GOVERNANCE_TAGS['risk_level']: risk_level,
 self.GOVERNANCE_TAGS['approval_status']: approval_status,
 self.GOVERNANCE_TAGS['fairness_status']: 'pass' if fairness_metrics.get('dp_ratio', 0) >= 0.8 else 'fail',
 }
 
 if tags:
 governance_tags.update(tags)
 
 # Model Version Tags Setup
 for key, value in governance_tags.items():
 self.client.set_model_version_tag(name, model_version, key, value)
 
 # Fairness (description) Addition
 description = f"""
## Governance Information

- Risk Level: {risk_level}
- Approval Status: {approval_status}
- DP Ratio: {fairness_metrics.get('dp_ratio', 'N/A')}
- EO Diff: {fairness_metrics.get('eo_diff', 'N/A')}
- 4/5 Rule: {'Pass' if fairness_metrics.get('dp_ratio', 0) >= 0.8 else 'Fail'}
"""
 
 self.client.update_model_version(
 name=name,
 version=model_version,
 description=description
 )
 
 return model_version
 
 def approve_model_version(
 self,
 name: str,
 version: str,
 approver: str,
 conditions: Optional[List[str]] = None,
 ) -> None:
 """Model Version Approval"""
 
 approval_date = datetime.now().isoformat()
 
 self.client.set_model_version_tag(
 name, version,
 self.GOVERNANCE_TAGS['approval_status'], 'approved'
 )
 self.client.set_model_version_tag(
 name, version,
 self.GOVERNANCE_TAGS['approved_by'], approver
 )
 self.client.set_model_version_tag(
 name, version,
 self.GOVERNANCE_TAGS['approval_date'], approval_date
 )
 
 if conditions:
 self.client.set_model_version_tag(
 name, version,
 f"{self.GOVERNANCE_PREFIX}.approval_conditions",
 json.dumps(conditions)
 )
 
 # (Staging -> Production)
 self.client.transition_model_version_stage(
 name=name,
 version=version,
 stage="Production",
 archive_existing_versions=True
 )
 
 def get_governance_status(
 self,
 name: str,
 version: Optional[str] = None,
 ) -> Dict:
 """Model Governance Status """
 
 if version is None:
 # Version 
 versions = self.client.search_model_versions(f"name='{name}'")
 if not versions:
 raise ValueError(f"Model {name} not found")
 version = max(v.version for v in versions)
 
 model_version = self.client.get_model_version(name, version)
 
 governance_info = {}
 for tag_name, tag_key in self.GOVERNANCE_TAGS.items():
 if tag_key in model_version.tags:
 governance_info[tag_name] = model_version.tags[tag_key]
 
 return {
 'model_name': name,
 'version': version,
 'stage': model_version.current_stage,
 'governance': governance_info,
 }
\end{lstlisting}


\subsection{CI/CD 파이프라인 통합된}
\label{sec:6.4.2}

AI 거버넌스 검사를 CI/CD Pipeline에 통합하여 
자동화된 품질 Gate를 구현한다.

\begin{lstlisting}[language=yaml, caption={GitHub Actions Governance Pipeline}, label={lst:github_actions}]
# .github/workflows/ai-governance.yml

name: AI Governance Checks

on:
  push:
    branches: [main, develop]
    paths:
      - 'models/'
      - 'src/'
  pull_request:
    branches: [main]

env:
  PYTHON_VERSION: '3.11'
  FAIRNESS_THRESHOLD: 0.8
  
jobs:
  fairness-check:
    name: Fairness Evaluation
    runs-on: ubuntu-latest
    
    steps:
      - uses: actions/checkout@v4
      
      - name: Set up Python
        uses: actions/setup-python@v5
        with:
          python-version: ${{ env.PYTHON_VERSION }}
      
      - name: Install dependencies
        run: |
          pip install -r requirements.txt
          pip install fairlearn==0.10.1 shap==0.44.0
      
      - name: Run Fairness Tests
        id: fairness
        run: |
          python -m pytest tests/test_fairness.py \
          --fairness-threshold=${{ env.FAIRNESS_THRESHOLD }} \
          --junitxml=fairness-results.xml
        continue-on-error: true
      
      - name: Upload Fairness Report
        uses: actions/upload-artifact@v4
        with:
          name: fairness-report
          path: reports/fairness_report.md
          retention-days: 1825 # 5years 
      
      - name: Check Fairness Gate
        if: steps.fairness.outcome == 'failure'
        run: |
          echo "::error::Fairness check failed. DP Ratio below threshold."
          exit 1

  explainability-check:
    name: Explainability Validation
    runs-on: ubuntu-latest
    needs: fairness-check
    
    steps:
      - uses: actions/checkout@v4
      
      - name: Set up Python
        uses: actions/setup-python@v5
        with:
          python-version: ${{ env.PYTHON_VERSION }}
      
      - name: Generate Model Card
        run: |
          python scripts/generate_model_card.py \
          --model-path models/latest \
          --output reports/model_card.json
      
      - name: Validate Model Card Completeness
        run: |
          python scripts/validate_model_card.py \
          --model-card reports/model_card.json \
          --schema schemas/model_card_schema.json
      
      - name: Upload Model Card
        uses: actions/upload-artifact@v4
        with:
          name: model-card
          path: reports/model_card.json
          retention-days: 1825
  
  governance-approval:
    name: Governance Gate
    runs-on: ubuntu-latest
    needs: [fairness-check, explainability-check]
    if: github.ref == 'refs/heads/main'
    environment: production
    
    steps:
      - name: Request Governance Approval
        run: |
          echo "Governance checks passed. Requesting approval..."
          # Approval Workflow Slack 
      
      - name: Wait for Approval
        # Approval (GitHub Environment Protection)
        run: echo "Approved for production deployment"
\end{lstlisting}

%----------------------------------------------------------
\section{모니터링 및 관측성 플랫폼}
\label{sec:6.5}

\subsection{Evidently AI를 활용한 모니터링}
\label{sec:6.5.1}

Evidently AI는 오픈소스 ML 모니터링 Tool로, 
데이터 드리프트, 모델 성능, Fairness을 추적할 수 있다.

\begin{lstlisting}[language=Python, caption={Evidently AI Monitoring Integration}, label={lst:evidently_monitoring}]
"""
Evidently AI Monitoring Setup

- Data Drift Detection
- Performance Monitoring
- Fairness Tracking
- Automated Report Generation
"""

import pandas as pd
import numpy as np
from datetime import datetime
from typing import Optional, Dict, List

from evidently import ColumnMapping
from evidently.report import Report
from evidently.metric_preset import (
 DataDriftPreset,
 DataQualityPreset,
 ClassificationPreset,
)
from evidently.metrics import (
 DataDriftTable,
 DatasetDriftMetric,
 ColumnDriftMetric,
 ClassificationQualityMetric,
 ClassificationClassBalance,
)
from evidently.test_suite import TestSuite
from evidently.test_preset import (
 DataDriftTestPreset,
 DataQualityTestPreset,
)
from evidently.tests import (
 TestColumnDrift,
 TestShareOfDriftedColumns,
 TestAccuracyScore,
)

class EvidentlyGovernanceMonitor:
 """
 Evidently AI based Governance Monitoring
 """
 
 def __init__(
 self,
 reference_data: pd.DataFrame,
 column_mapping: Optional[ColumnMapping] = None,
 ):
 """
 Parameters
 ----------
 reference_data : pd.DataFrame
 Data (Learning Data Data)
 column_mapping : ColumnMapping
 (target, prediction, numerical, categorical )
 """
 self.reference_data = reference_data
 self.column_mapping = column_mapping or ColumnMapping()
 
 def create_drift_report(
 self,
 current_data: pd.DataFrame,
 output_path: Optional[str] = None,
 ) -> Report:
 """Data Drift Report Generation"""
 
 report = Report(metrics=[
 DataDriftPreset(),
 DatasetDriftMetric(),
 ])
 
 report.run(
 reference_data=self.reference_data,
 current_data=current_data,
 column_mapping=self.column_mapping,
 )
 
 if output_path:
 report.save_html(output_path)
 
 return report
 
 def create_performance_report(
 self,
 current_data: pd.DataFrame,
 output_path: Optional[str] = None,
 ) -> Report:
 """Performance Monitoring Report Generation"""
 
 report = Report(metrics=[
 ClassificationPreset(),
 ClassificationClassBalance(),
 ])
 
 report.run(
 reference_data=self.reference_data,
 current_data=current_data,
 column_mapping=self.column_mapping,
 )
 
 if output_path:
 report.save_html(output_path)
 
 return report
 
 def run_drift_tests(
 self,
 current_data: pd.DataFrame,
 drift_threshold: float = 0.1,
 ) -> Dict:
 """Drift """
 
 test_suite = TestSuite(tests=[
 TestShareOfDriftedColumns(lt=drift_threshold),
 DataDriftTestPreset(),
 ])
 
 test_suite.run(
 reference_data=self.reference_data,
 current_data=current_data,
 column_mapping=self.column_mapping,
 )
 
 # results = {
 'passed': test_suite.as_dict()['summary']['all_passed'],
 'total_tests': test_suite.as_dict()['summary']['total_tests'],
 'failed_tests': test_suite.as_dict()['summary']['failed_tests'],
 'details': test_suite.as_dict()['tests'],
 }
 
 return results
 
 def monitor_fairness_over_time(
 self,
 data_batches: List[pd.DataFrame],
 sensitive_column: str,
 prediction_column: str = 'prediction',
 target_column: str = 'target',
 ) -> pd.DataFrame:
 """
 Fairness Tracking
 
 Data Fairness Calculation
 """
 
 results = []
 
 for i, batch in enumerate(data_batches):
 groups = batch[sensitive_column].unique()
 
 # Calculation
 selection_rates = {}
 for group in groups:
 mask = batch[sensitive_column] == group
 selection_rates[group] = batch.loc[mask, prediction_column].mean()
 
 # DP Ratio Calculation
 if len(selection_rates) >= 2:
 rates = list(selection_rates.values())
 dp_ratio = min(rates) / max(rates) if max(rates) > 0 else 1.0
 else:
 dp_ratio = 1.0
 
 results.append({
 'batch': i,
 'dp_ratio': dp_ratio,
 {f'selection_rate_{g}': r for g, r in selection_rates.items()},
 'passes_4_5_rule': dp_ratio >= 0.8,
 })
 
 return pd.DataFrame(results)
 
 def generate_governance_dashboard_data(
 self,
 current_data: pd.DataFrame,
 ) -> Dict:
 """Governance Data Generation"""
 
 # Drift Analysis
 drift_report = self.create_drift_report(current_data)
 drift_results = drift_report.as_dict()
 
 # Performance Analysis
 perf_report = self.create_performance_report(current_data)
 perf_results = perf_report.as_dict()
 
 # Data 
 dashboard_data = {
 'timestamp': datetime.now().isoformat(),
 'data_quality': {
 'total_rows': len(current_data),
 'missing_rate': current_data.isnull().mean().mean(),
 },
 'drift': {
 'dataset_drift': drift_results.get('metrics', [{}])[0].get('result', {}),
 'drifted_columns': [], # Drift },
 'performance': {
 'metrics': perf_results.get('metrics', [{}])[0].get('result', {}),
 },
 'alerts': [],
 }
 
 # Generation
 if dashboard_data['drift'].get('dataset_drift', {}).get('drift_share', 0) > 0.2:
 dashboard_data['alerts'].append({
 'severity': 'warning',
 'message': 'Data Drift 20% .',
 })
 
 return dashboard_data
\end{lstlisting}


%----------------------------------------------------------
\section{문서화 자동화 도구}
\label{sec:6.6}

\subsection{Model Card Toolkit}
\label{sec:6.6.1}

Google의 모델 Card Toolkit을 활용하여 
표준화된 모델 카드를 자동 생성할 수 있다.

\begin{lstlisting}[language=Python, caption={Integrated Documentation Automation System}, label={lst:documentation_automation}]
"""
AI Governance Documentation Automation

- Automated Model Card Generation
- Fairness Report Automation
- Compliance Document Package
"""

from dataclasses import dataclass, field
from typing import Dict, List, Optional, Any
from datetime import datetime
import json
import os

@dataclass
class GovernanceDocumentPackage:
 """
 AI Governance Document 15 EU AI Act Annex IV Compliance Document Generation
 """
 
 # System information
 system_name: str
 system_id: str
 version: str
 developer: str
 
 # Classification
 risk_level: str # 'high', 'medium', 'low'
 high_risk_sector: Optional[str] = None
 
 # Model information
 model_type: str = ""
 intended_use: List[str] = field(default_factory=list)
 out_of_scope_use: List[str] = field(default_factory=list)
 
 # Data information
 training_data_source: str = ""
 training_data_size: int = 0
 training_data_period: str = ""
 
 # Performance Fairness
 performance_metrics: Dict[str, float] = field(default_factory=dict)
 fairness_metrics: Dict[str, float] = field(default_factory=dict)
 subgroup_performance: Dict[str, Dict[str, float]] = field(default_factory=dict)
 
 # 
 xai_method: str = ""
 feature_importance: Dict[str, float] = field(default_factory=dict)
 
 # human_oversight_level: str = ""
 override_rate: float = 0.0
 
 # identified_risks: List[Dict] = field(default_factory=list)
 mitigation_measures: List[str] = field(default_factory=list)
 
 # Data
 creation_date: datetime = field(default_factory=datetime.now)
 retention_years: int = 5
 
 def generate_model_card(self) -> str:
 """Model Generation ( )"""
 
 card = f"""# Model : {self.system_name}

## 1. Model 

| Items | |
|------|------|
| System | {self.system_name} |
| System ID | {self.system_id} |
| Version | {self.version} |
| | {self.developer} |
| Model | {self.model_type} |
| Level | {self.risk_level.upper()} |
"""
 
 if self.high_risk_sector:
 card += f"| | {self.high_risk_sector} |\n"
 
 card += f"""
## 2. ### Case
"""
 for use in self.intended_use:
 card += f"- {use}\n"
 
 card += "\n### ( )\n"
 for use in self.out_of_scope_use:
 card += f"- {use}\n"
 
 card += f"""

## 3. Learning Data

| Items | |
|------|------|
| Data | {self.training_data_source} |
| Data | {self.training_data_size:,} |
| Period | {self.training_data_period} |

## 4. Performance

| | |
|--------|-----|
"""
 for metric, value in self.performance_metrics.items():
 card += f"| {metric} | {value:.4f} |\n"
 
 card += f"""

## 5. Fairness Assessment

| | | |
|--------|-----|----------|
"""
 for metric, value in self.fairness_metrics.items():
 passed = '[PASS]' if metric == 'dp_ratio' and value >= 0.8 else '[PASS]' if metric == 'eo_diff' and value <= 0.1 else '[WARN]'
 card += f"| {metric} | {value:.4f} | {passed} |\n"
 
 if self.subgroup_performance:
 card += "\n### Performance\n\n"
 card += "| | Accuracy | Precision | Recall |\n"
 card += "|------|----------|-----------|--------|\n"
 for group, metrics in self.subgroup_performance.items():
 card += f"| {group} | {metrics.get('accuracy', 0):.4f} | {metrics.get('precision', 0):.4f} | {metrics.get('recall', 0):.4f} |\n"
 
 card += f"""

## 6. 

-XAI : {self.xai_method}

### Major ( 10 )

| | | |
|------|------|--------|
"""
 sorted_importance = sorted(
 self.feature_importance.items(),
 key=lambda x: abs(x[1]),
 reverse=True
 )[:10]
 for i, (feat, imp) in enumerate(sorted_importance, 1):
 card += f"| {i} | {feat} | {imp:.4f} |\n"
 
 card += f"""

## 7. | Items | |
|------|------|
| Level | {self.human_oversight_level} |
| AI | {self.override_rate:.2%} |

## 8. ### """
 for risk in self.identified_risks:
 card += f"-{risk.get('name', 'N/A')}** ( : {risk.get('severity', 'N/A')}): {risk.get('description', '')}\n"
 
 card += "\n### \n"
 for measure in self.mitigation_measures:
 card += f"- {measure}\n"
 
 card += f"""

---

## Document information

| Items | |
|------|------|
| days | {self.creation_date.strftime('%Y-%m-%d')} |
| | {self.creation_date.year + self.retention_years}years ({self.retention_years}years) |

* Document AI 15 .*
"""
 
 return card
 
 def generate_compliance_checklist(self) -> str:
 """Regulation Compliance Checklist Generation"""
 
 checklist = f"""# Regulation Compliance Checklist

System: {self.system_name} ({self.system_id}) 
Assessmentdays: {datetime.now().strftime('%Y-%m-%d')}

---

## AI Compliance

| | | Compliance | |
|------|----------|----------|------|
| 10 | AI Classification | {'[PASS]' if self.risk_level == 'high' and self.high_risk_sector else '[PASS]' if self.risk_level != 'high' else ''} | Classification Document |
| 12 | | {'[PASS]' if self.fairness_metrics.get('dp_ratio', 0) >= 0.8 else '[FAIL]'} | Fairness Assessment Report |
| 13 | | | |
| 14 | | {'[PASS]' if self.human_oversight_level in ['HITL', 'HOTL'] else '[WARN]'} | Document |
| 15 | (5years) | [PASS] | Document |
| 23 | | {'[PASS]' if self.xai_method else '[FAIL]'} | Model , XAI Report |

---

## EU AI Act Compliance ( )

| | | Compliance |
|------|----------|----------|
| Art. 9 | Management System | {'[PASS]' if self.identified_risks else '[FAIL]'} |
| Art. 10 | Data Governance | |
| Art. 11 | Document | [PASS] |
| Art. 13 | | [PASS] |
| Art. 14 | | {'[PASS]' if self.human_oversight_level in ['HITL', 'HOTL'] else '[WARN]'} |

---

Check: ________________ 
days: ________________ 
 : ________________
"""
 
 return checklist
 
 def export_package(self, output_dir: str) -> List[str]:
 """ Document """
 
 os.makedirs(output_dir, exist_ok=True)
 
 files_created = []
 
 # 1. Model 
 model_card_path = os.path.join(output_dir, "model_card.md")
 with open(model_card_path, "w", encoding="utf-8") as f:
 f.write(self.generate_model_card())
 files_created.append(model_card_path)
 
 # 2. Compliance Checklist
 checklist_path = os.path.join(output_dir, "compliance_checklist.md")
 with open(checklist_path, "w", encoding="utf-8") as f:
 f.write(self.generate_compliance_checklist())
 files_created.append(checklist_path)
 
 # 3. Data JSON
 metadata = {
 "system_name": self.system_name,
 "system_id": self.system_id,
 "version": self.version,
 "risk_level": self.risk_level,
 "high_risk_sector": self.high_risk_sector,
 "fairness_metrics": self.fairness_metrics,
 "performance_metrics": self.performance_metrics,
 "human_oversight_level": self.human_oversight_level,
 "creation_date": self.creation_date.isoformat(),
 "retention_until": (self.creation_date.year + self.retention_years),
 }
 
 metadata_path = os.path.join(output_dir, "metadata.json")
 with open(metadata_path, "w", encoding="utf-8") as f:
 json.dump(metadata, f, ensure_ascii=False, indent=2)
 files_created.append(metadata_path)
 
 return files_created
\end{lstlisting}


%----------------------------------------------------------
\subsection*{제6장 요약}

본 장에서는 AI 거버넌스를 기술적으로 Support하는 도구 생태계를 다루었다. 
핵심 내용을 요약하면 다음과 같다.

\begin{enumerate}
 \item \textbf{기술 스택 아키텍처}: 
 5개 계층(MLOps 인프라 -> 모니터링 -> 평가 도구 -> 정책 관리 -> 대시보드)으로 
 구성된 AI 거버넌스 기술 스택을 제시하였다.
 
 \item \textbf{공정성 도구}: 
 Fairlearn 0.14.0의 고급 기능(다중 속성 분석, 교차 분석, 완화 알고리즘 비교)과 
 AIF360과의 비교를 다루었다.
 
 \item \textbf{설명가능한 AI(설명가능한 AI(XAI)) 도구}: 
 SHAP 0.46.0의 거버넌스 통합 활용법과 
 CaptumAdvanced Deep Learning with 모델 설명 방법을 제시하였다.
 
 \item \textbf{MLOps 통합된}: 
 MLflow와 거버넌스 정보 Integration, GitHub Actions를 활용한 
 CI/CD 파이프라인 내 거버넌스 Gate 구현을 다루었다.
 
 \item \textbf{모니터링}: 
 Evidently AI를 활용한 드리프트 감지, 성능 모니터링, 
 공정성 추적 방법을 제시하였다.
 
 \item \textbf{문서화 자동화}: 
 시행령 제15조 및 EU AI Act Annex IV 준수를 위한 
 문서 패키지 자동 생성 Tool를 제시하였다.
\end{enumerate}

이러한 Tool들을 조직의 기존 ML 인프라에 통합함으로써, 수작업 부담을 줄이고 일관된 거버넌스 실행을 보장할 수 있다. 이어지는 제7장에서는 이러한 거버넌스 Tool와 프레임워크가 실제 한국 기업의 현장에서 어떻게 적용되고 있는지 구체적인 사례를 통해 살펴본다.
